% ---------------------------------------------------------------------------
% Related Work (Section 2) - drafted for M2 (keyword spotting, recurrence,
% attention, augmentation). M4 adds the self-supervised speech models and
% African/low-resource speech subsection. Citation keys match report/references.bib.
% ---------------------------------------------------------------------------
\section{Related Work}
\label{sec:related}

\subsection{From HMMs to neural keyword spotting}
Recognising a word from a small vocabulary is one of the oldest speech tasks. Classic systems pair a compact
spectral representation, the mel-frequency cepstrum, which outperformed other parametric representations for
word recognition~\cite{davis1980comparison}, with one hidden Markov model per word, trained by Baum--Welch
re-estimation and decoded by maximum likelihood~\cite{rabiner1989tutorial}. Our A1 and A2 baselines follow this
tradition. A1 keeps only order-free statistics of the cepstral frames, and A2 is the per-word left-to-right HMM.

The task became a standard benchmark with the Speech Commands dataset, which also proposed a reproducible way to
report keyword-spotting accuracy~\cite{warden2018speech}. Convolutional networks then replaced HMM and DNN
pipelines for small-footprint keyword spotting. Under fixed budgets on multiplications and parameters, they
reduced the false-reject rate by 27--44\% relative to a DNN~\cite{sainath2015convolutional}. Temporal
convolutions over the frames, arranged in a compact residual network (TC-ResNet), made keyword spotting fast
enough for real-time use on phones (a reported 385$\times$ speed-up) while surpassing the accuracy of the previous
state of the art on Speech Commands~\cite{choi2019temporal}. This is the design behind our A4.

\subsection{Recurrence and attention}
Recurrent networks model a sequence by carrying a hidden state from frame to frame. Long short-term memory (LSTM)
units keep the error signal constant through gated memory cells, and can learn dependencies across more than
1{,}000 time steps~\cite{hochreiter1997long}. Stacked, bidirectional LSTMs brought this to acoustic modelling,
reaching a 17.7\% phoneme error rate on TIMIT~\cite{graves2013speech}. Attention replaces a single summary
vector with a learned, soft search over positions~\cite{bahdanau2015neural}. For speech commands, de Andrade et
al.\ combined convolution, recurrence and attention in one model, reaching 94.1\% and 94.5\% accuracy on Speech
Commands V1 and V2 with only 202K parameters. They also used the attention weights to show which parts of the
audio drove each decision~\cite{deandrade2018neural}.

Our A3 belongs to this line. It differs in separating the contributions: A3 is recurrent by default, the
convolutional front-end is tested as an ablation, and A4 is purely convolutional.

Whether recurrence is needed at all is itself an open question. A systematic comparison across sequence-modelling
benchmarks found that a simple convolutional architecture outperforms canonical recurrent networks such as LSTMs,
while also showing a longer effective memory~\cite{bai2018empirical}. Our A3--A4 comparison tests whether this
holds for sequences as short as a single spoken word, with only a few hundred training examples per class.

\subsection{Augmentation for small speech datasets}
SpecAugment showed that simple transformations applied directly to filter-bank features (time warping, and
masking blocks of frequency channels and of time steps) are enough to reach state-of-the-art results on large
speech-recognition benchmarks~\cite{park2019specaugment}. Our setting is the opposite extreme: about 245 training
clips per word, and competition rules that forbid external data such as noise corpora. Synthetic, feature-level
augmentation is therefore the only augmentation available to us, and its value is tested directly in round R2.

\subsection{Self-supervised speech models and African languages}
[M4: wav2vec~2.0 / XLS-R / MMS pretraining and fine-tuning; Common Voice; keyword spotting and speech technology
for African and other low-resource languages; the gap this study addresses.]

\subsection{How the literature shaped our study}
Three lessons from this work shaped the design:
\begin{enumerate}
  \item Keyword spotting has been solved with very different ways of modelling time: order-free features,
        Markov states, convolution, recurrence and self-attention. We therefore compare one representative of
        each on the same data, split and metrics, rather than tuning a single family.
  \item Attention-based recognisers can show where in the audio they look~\cite{deandrade2018neural}. We use this
        to test our central hypothesis directly: words whose sounds differ mainly in order (\textit{tisa} and
        \textit{sita}) should separate order-aware from order-free models.
  \item Augmentation such as SpecAugment is most useful when data is scarce~\cite{park2019specaugment}. With 245
        clips per word, we treat it as a first-class experimental factor rather than a fixed default.
\end{enumerate}
